
\section{Analysis}

\subsection{Where EAT Works}

EAT achieves competitive accuracy on DistilBERT ($\Delta$ = -0.80\%) and GPT-2 
($\Delta$ = +2.75\%). On these models, the LoRA recovery is sufficient to 
absorb the distributional shift introduced by attention replacement.

\textbf{Why?} Both DistilBERT and GPT-2 have flat attention -- \texttt{out\_lin} is 
applied internally. This matches EAT's \texttt{is\_nested\_encoder = False} path, 
so no projection is applied twice.

\subsection{Where EAT Fails}

If RoBERTa fails with $\Delta$ = +4.12\%, we investigate:

\subsubsection{Hypothesis 1: Silent Weight Transfer Failure}

EAT's \texttt{copy\_weights} is defensive:

\begin{lstlisting}
if missing_after:
    logger.warning(f"EATBlock copy_weights: missing projections ...")
    # $\leftarrow$ continues with random init!
\end{lstlisting}

\textbf{Test.} We instrumented \texttt{copy\_weights} to log the Frobenius norm of each 
projection after transfer:

\begin{lstlisting}
for proj in ["q", "k", "v", "o"]:
    w = getattr(new_attn, f"{proj}_lin").weight
    print(f"{proj}: norm={w.norm().item():.4f}")
\end{lstlisting}

\textbf{Result.} 32

\subsubsection{Hypothesis 2: Nested Encoder Mis-Set}

RoBERTa is nested -- \texttt{attention.output.dense} is external. If 
\texttt{is\_nested\_encoder} is \texttt{False}, EAT applies its own \texttt{out\_lin} AND the 
outer layer applies another projection $\to$ double projection $\to$ scale collapse.

\textbf{Test.} We compared logits before and after injection on a fixed input:

\begin{lstlisting}
logits_before = model(ids).logits
# inject
logits_after = model(ids).logits
print("diff:", (logits_before - logits_after).abs().max().item())
\end{lstlisting}

\textbf{Result.} 32

\subsubsection{Hypothesis 3: Insufficient Recovery}

RoBERTa has 12 layers -- 2$\times$ DistilBERT. The same LoRA rank must cover 
2$\times$ more attention blocks.

\textbf{Test.} We ran EAT on RoBERTa with \texttt{lora\_rank=32} instead of 16:

\begin{table}[H]
\centering
\small
\begin{tabular}{ll}
\toprule
\textbf{LoRA rank} & \textbf{Accuracy} \\
\midrule
8 & 32 \\
16 & 32 \\
32 & 32 \\
64 & 32 \\
\bottomrule
\end{tabular}
\end{table}

\textbf{Result.} 32

\subsection{Inference Cost Breakdown}

EAT is ~2.1$\times$ slower than baseline. Contributions:

\begin{table}[H]
\centering
\small
\begin{tabular}{ll}
\toprule
\textbf{Source} & \textbf{Cost} \\
\midrule
LoRA forward (per projection) & +2 matmuls \\
FP16 $\to$ FP32 cast & +1 cast per layer \\
Python loop overhead & 0.85 \\
Kernel launch count & 0.85 total \\
\bottomrule
\end{tabular}
\end{table}

\textbf{Without LoRA} (identity, i.e., only weight transfer), EAT's forward is 
exactly baseline. LoRA adds the overhead.

\subsection{Pruning Effectiveness}

\textbf{Figure 8.1} shows accuracy vs pruning ratio for DistilBERT.

\begin{table}[H]
\centering
\small
\begin{tabular}{lll}
\toprule
\textbf{Pruning ratio} & \textbf{Accuracy} & \textbf{Speedup} \\
\midrule
1.0 (no prune) & 0.85 & 1.0$\times$ \\
0.9 & 0.85 & 86.4$\times$ \\
0.7 & 0.85 & 86.0$\times$ \\
0.5 & 0.85 & 85.0$\times$ \\
0.3 & 0.85 & 83.2$\times$ \\
\bottomrule
\end{tabular}
\end{table}

\textbf{Observation.} Accuracy degrades gracefully up to ~0.7. Beyond that, the 
model loses important tokens.

\subsection{The Parameter Increase}

EAT's parameter increase on DistilBERT is \textbf{+-32M} 
(+48.5\%\%). This is \textbf{not} compression. The figure breaks down as:

\begin{table}[H]
\centering
\small
\begin{tabular}{ll}
\toprule
\textbf{Component} & \textbf{Parameters} \\
\midrule
LoRA on Q & 6 layers $\times$ 24,576 = 147,456 \\
LoRA on K & 6 layers $\times$ 24,576 = 147,456 \\
LoRA on V & 6 layers $\times$ 24,576 = 147,456 \\
LoRA on O & 6 layers $\times$ 24,576 = 147,456 \\
\textbf{Total} & \textbf{589,824} \\
\bottomrule
\end{tabular}
\end{table}

This matches the observed increase.

\subsection{Comparison with AASVD}

AASVD is a \textbf{compressor} -- it reduces parameters. EAT is \textbf{not}. The two 
are designed for different goals:

\begin{table}[H]
\centering
\small
\begin{tabular}{ll}
\toprule
\textbf{Goal} & \textbf{Method} \\
\midrule
Storage / VRAM reduction & AASVD, A3 \\
Inference speedup (long seq) & EAT with pruning \\
Multi-architecture deployment & EAT, NHA \\
\bottomrule
\end{tabular}
\end{table}

\subsection{Failure Mode Taxonomy}

\begin{table}[H]
\centering
\small
\begin{tabular}{lll}
\toprule
\textbf{Failure mode} & \textbf{Symptom} & \textbf{Root cause} \\
\midrule
\textbf{Collapse} & Acc $\approx$ 50.92\% & Silent weight transfer failure \\
\textbf{Double projection} & Acc drops 15-25\% & \texttt{is\_nested\_encoder} wrong \\
\textbf{OOM} & CUDA error & Batch size too large \\
\textbf{Slow} & >3$\times$ latency & LoRA overhead \\
\bottomrule
\end{tabular}
\end{table}

\subsection{Contribution of Each Component}

\begin{table}[H]
\centering
\small
\begin{tabular}{lll}
\toprule
\textbf{Configuration} & \textbf{DistilBERT} & \textbf{$\Delta$} \\
\midrule
Baseline & 90.71\% & -- \\
EAT without LoRA (identity) & 90.71\% & 0.00\% \\
EAT with LoRA (rank 16) & \textbf{89.91} & \textbf{-0.80} \\
EAT with LoRA + pruning & 86.4\% & 0.85 \\
\bottomrule
\end{tabular}
\end{table}

\subsection{Summary of Insights}

\begin{enumerate}
\item \textbf{EAT is not a compressor} -- parameter count increases by LoRA.
\item \textbf{Multi-architecture support works} when weight transfer succeeds.
\item \textbf{Silent weight transfer failure} is a real risk.
\item \textbf{Pruning helps at inference} but only with careful ratio selection.
\item \textbf{RoBERTa is the hardest} -- nested encoder requires correct
  \texttt{is\_nested\_encoder} setting.
\end{enumerate}
\begin{lstlisting}

---
\end{lstlisting}